\section{Introduction}
\label{sec:intro}

World models learn to predict future observations and can use those predictions for planning. Good
video prediction, however, does not by itself show that a model has learned the underlying dynamics
of the system it observes. That distinction matters most when we ask the model to predict far beyond
the trajectories it saw during training.

A common way to look for physical structure is to fit a probe from the hidden state to a physical
variable. Probe accuracy tells us what information a latent representation retains, but not whether
the model's transition uses that information. Our own baseline makes this problem concrete. Six
randomly initialized DreamerV3 models contain polynomial functions of their latent state that
correlate with pendulum energy at up to $0.908$. A strong correlation can therefore appear even
before learning.

\begin{figure}[t]
\centering
\includegraphics[width=\textwidth]{figures/fig1_three_claims.pdf}
\caption{Each point represents one model. \textbf{(a)} The three trained conservative models recover
energy correlations between $0.967$ and $0.975$, while six untrained models span $0.17$--$0.91$.
\textbf{(b)} The recovered scalar changes very little within observation-conditioned trajectories
for conservative models and much more for damped models. \textbf{(c)} Rollout error as a function of
projection strength $\alpha$, relative to no edit. Enforcing the recovered constraint lowers error;
enforcing a matched random polynomial constraint usually raises it; applying the recovered
constraint to a damped model changes little.}
\label{fig:three}
\end{figure}

We therefore ask whether a trained world model contains a scalar that its own transition
approximately preserves, and whether failures to preserve that scalar contribute to prediction
error. We study a frozen DreamerV3 \citep{hafner2023dreamerv3} trained only on $64\times64$ video of
a Gymnasium pendulum, with no physical labels and no actor or critic.

Independently trained models recover nearly identical energy-like invariants. Decodability alone
does not explain this: randomly initialized Dreamers also produce strongly energy-correlated
readouts, while matched models trained on damped dynamics contain no comparably conserved scalar.
The recovered invariant then begins to drift during autonomous imagination. Correcting that drift
improves 50-step predictions in every conservative model, whereas matched random corrections usually
make them worse. Dreamer has learned a physical constraint that its own rollout dynamics fail to
preserve.
